\documentclass[11pt]{article}
% ===========================================================================
%  Shared preamble for the AP Calculus AB exam-preparation worksheets.
%  Input right after \documentclass. Files beginning with "_" are not compiled
%  on their own. Navy and gold, to match the course's lesson pages.
% ===========================================================================
\usepackage[letterpaper,top=2.3cm,bottom=2.2cm,left=1.7cm,right=1.7cm,headheight=28pt,footskip=30pt]{geometry}
\usepackage{amsmath,amssymb}
\usepackage[T1]{fontenc}
\usepackage[utf8]{inputenc}
\usepackage{lmodern}
\usepackage{xcolor}
\usepackage[most]{tcolorbox}
\usepackage{enumitem}
\usepackage{fancyhdr}
\usepackage{array}
\usepackage{tikz}
\usepackage{pgfplots}
\pgfplotsset{compat=1.18}

\definecolor{navy}{HTML}{161E54}
\definecolor{ink}{HTML}{27357E}
\definecolor{gold}{HTML}{B98A00}
\definecolor{goldbg}{HTML}{FFFDF3}
\definecolor{inkbg}{HTML}{F6F7FD}
\definecolor{rule}{HTML}{C9CFF0}
\definecolor{okgreen}{HTML}{1B7A44}
\definecolor{okbg}{HTML}{F1F8F3}

\pagestyle{fancy}
\fancyhf{}
\fancyhead[L]{\footnotesize Name: \underline{\hspace{4cm}}}
\fancyhead[C]{\textbf{\color{navy}Integration Academy}\\[-2pt]{\scriptsize Calculus AB \textperiodcentered{} AP\textsuperscript{\textregistered} Exam Preparation}}
\fancyfoot[L]{\tiny\color{black!55}\parbox[t]{0.78\textwidth}{AP\textsuperscript{\textregistered} is a registered trademark of the College Board, which was not involved in the production of, and does not endorse, this material. \textcopyright{} Integration Academy.}}
\fancyfoot[R]{\scriptsize\thepage}
\renewcommand{\headrulewidth}{1.2pt}
\renewcommand{\headrule}{\hbox to\headwidth{\color{gold}\leaders\hrule height \headrulewidth\hfill}}
\renewcommand{\footrulewidth}{0pt}
\setlength{\parindent}{0pt}
\setlength{\parskip}{4pt}

% Title block: \wstitle{code}{title}{unit line}
\newcommand{\wstitle}[3]{\begin{center}{\small\bfseries\color{gold}WORKSHEET #1}\\[3pt]{\LARGE\bfseries\color{navy}#2}\\[3pt]{\small\color{black!60}#3}\end{center}\vspace{2pt}}
% Section heading
\newcommand{\sech}[1]{\par\vspace{10pt}{\large\bfseries\color{navy}#1}\par\nobreak\vspace{2pt}{\color{rule}\hrule height0.8pt}\vspace{7pt}}
% A level tag printed beside a group of questions
\newcommand{\level}[1]{\par\vspace{6pt}{\footnotesize\bfseries\color{gold}\MakeUppercase{#1}}\par\nobreak\vspace{2pt}}

\newtcolorbox{ideabox}{colframe=ink,colback=inkbg,coltitle=white,colbacktitle=ink,fonttitle=\bfseries\small,title={Key ideas},arc=2pt,boxrule=0.8pt,left=8pt,right=8pt}
\newtcolorbox{example}[1]{colframe=rule,colback=white,coltitle=navy,colbacktitle=inkbg,fonttitle=\bfseries\small,title={#1},arc=2pt,boxrule=0.8pt,borderline west={3pt}{0pt}{ink},left=10pt,right=8pt}
\newtcolorbox{frbox}[1]{colframe=gold,colback=goldbg,coltitle=white,colbacktitle=gold,fonttitle=\bfseries\small,title={#1},arc=2pt,boxrule=0.8pt,breakable,left=8pt,right=8pt}
\newtcolorbox{solbox}{colframe=okgreen,colback=okbg,coltitle=white,colbacktitle=okgreen,fonttitle=\bfseries\small,title={Worked solutions},arc=2pt,boxrule=0.8pt,breakable,left=8pt,right=8pt}
\newtcolorbox{scorebox}[1]{colframe=navy,colback=white,coltitle=white,colbacktitle=navy,fonttitle=\bfseries\small,title={#1},arc=2pt,boxrule=0.8pt,breakable,left=8pt,right=8pt}

\newcommand{\soln}{\par\smallskip\textbf{\color{ink}Solution.}\ }
\newcommand{\result}[1]{\par\smallskip\textbf{#1}}
% blank working space of a given height
\newcommand{\work}[1]{\par\vspace{#1}}
% a point value in the margin of a scoring guide
\newcommand{\pts}[1]{\hfill{\footnotesize\bfseries\color{gold}[#1]}}

% Coordinate grid for graph questions: \begin{apgraph}{xmin}{xmax}{ymin}{ymax} ... \end{apgraph}
\newenvironment{apgraph}[4]{\begin{center}\begin{tikzpicture}\begin{axis}[
  axis lines=middle, xmin=#1, xmax=#2, ymin=#3, ymax=#4,
  grid=both, grid style={black!12}, xtick distance=1, ytick distance=1,
  ticklabel style={font=\scriptsize}, xlabel={$x$}, ylabel={$y$},
  width=10cm, height=7.2cm, samples=80, clip=false,
  every axis plot/.append style={line width=1.1pt, color=ink}]}{\end{axis}\end{tikzpicture}\end{center}}
\newcommand{\closeddot}[2]{\addplot[only marks, mark=*, mark size=2.2pt, color=ink] coordinates {(#1,#2)};}
\newcommand{\opendot}[2]{\addplot[only marks, mark=*, mark size=2.2pt, color=ink, mark options={fill=white}] coordinates {(#1,#2)};}

% Multiple-choice options: four across, or stacked when the options are long.
\newcommand{\opts}[4]{\par\vspace{3pt}\begin{tabular}{@{}p{0.235\linewidth}p{0.235\linewidth}p{0.235\linewidth}p{0.235\linewidth}@{}}(A)\ #1 & (B)\ #2 & (C)\ #3 & (D)\ #4\end{tabular}\par\vspace{9pt}}
\newcommand{\optsv}[4]{\par\vspace{3pt}\begin{tabular}{@{}l@{\ \ }p{0.9\linewidth}@{}}(A) & #1\\(B) & #2\\(C) & #3\\(D) & #4\end{tabular}\par\vspace{9pt}}

\fancyhead[R]{\footnotesize\color{black!60}Worksheet 0.2}
\begin{document}

\wstitle{0.2}{Trigonometry Without Hesitation}{Unit 0: Readiness}

\begin{ideabox}
\begin{itemize}[leftmargin=*,itemsep=2pt,topsep=2pt]
  \item Calculus uses \textbf{radians}: $\pi$ radians is $180^\circ$.
  \item For an angle outside the first quadrant, use its \textbf{reference angle} and attach the sign: sine is positive in quadrants I and II, cosine in I and IV, tangent in I and III.
  \item $\sin^2x+\cos^2x=1$, \quad $\sin2x=2\sin x\cos x$, \quad $\cos2x=1-2\sin^2x=2\cos^2x-1$.
  \item To solve an equation, move everything to one side and \textbf{factor}. Never divide by $\sin x$ or $\cos x$.
\end{itemize}
\end{ideabox}

\sech{Worked examples}
\begin{example}{Example 1 \textperiodcentered{} Exact values}
Find $\cos\dfrac{3\pi}{4}$, $\sin\dfrac{11\pi}{6}$ and $\tan\dfrac{7\pi}{6}$.\soln Reference angles $\dfrac\pi4$, $\dfrac\pi6$, $\dfrac\pi6$, in quadrants II, IV and III.\result{$-\dfrac{\sqrt2}{2}$, \ $-\dfrac12$ \ and \ $\dfrac{1}{\sqrt3}$.}
\end{example}

\begin{example}{Example 2 \textperiodcentered{} A linear equation}
Solve $2\cos x+1=0$ for $0\le x<2\pi$.\soln $\cos x=-\dfrac12$. The reference angle is $\dfrac\pi3$ and cosine is negative in quadrants II and III.\result{$x=\dfrac{2\pi}{3}$ and $x=\dfrac{4\pi}{3}$.}
\end{example}

\begin{example}{Example 3 \textperiodcentered{} Using an identity first}
Solve $\cos2x+\sin x=0$ for $0\le x<2\pi$.\soln $1-2\sin^2x+\sin x=0$, so $2\sin^2x-\sin x-1=(2\sin x+1)(\sin x-1)=0$. $\sin x=1$ gives $\dfrac\pi2$; $\sin x=-\dfrac12$ gives $\dfrac{7\pi}{6}$ and $\dfrac{11\pi}{6}$.\result{$x=\dfrac\pi2,\ \dfrac{7\pi}{6},\ \dfrac{11\pi}{6}$.}
\end{example}

\newpage
\sech{Practice}

\level{Core}
\begin{enumerate}[leftmargin=*,itemsep=2pt]
  \item Find $\sin\dfrac{3\pi}{2}$ and $\cos\dfrac{2\pi}{3}$.\work{2.3cm}
  \item Convert $\dfrac{5\pi}{4}$ radians to degrees.\work{2.3cm}
  \item Solve $\sin x=\dfrac{\sqrt3}{2}$ for $0\le x<2\pi$.\work{2.3cm}
  \item Solve $\tan x=-1$ for $0\le x<2\pi$.\work{2.3cm}
\end{enumerate}

\level{Exam level}
\begin{enumerate}[leftmargin=*,itemsep=2pt,start=5]
  \item Solve $4\cos^2x=3$ for $0\le x<2\pi$.\work{2.8cm}
  \item Given $\cos\theta=-\dfrac23$ with $\theta$ in quadrant II, find $\sin\theta$ and $\tan\theta$.\work{2.8cm}
  \item Show that $\tan x+\cot x=\sec x\csc x$.\work{2.8cm}
  \item Find $\arcsin\dfrac{\sqrt3}{2}$, $\arctan\!\left(-\sqrt3\right)$ and $\arccos0$.\work{2.8cm}
\end{enumerate}

\begin{frbox}{Question 9 \textperiodcentered{} Free-response practice}
Let $g(x)=\cos2x+3\cos x+2$ for $0\le x<2\pi$.
\begin{enumerate}[label=(\alph*),leftmargin=*,itemsep=2pt]
  \item Find the exact value of $g\!\left(\dfrac\pi3\right)$.\work{2.6cm}
  \item Write $g(x)$ in terms of $\cos x$ only.\work{2.6cm}
  \item Find all values of $x$ in the interval at which $g(x)=0$.\work{2.6cm}
\end{enumerate}
\end{frbox}

\newpage
\begin{solbox}
\begin{enumerate}[leftmargin=*,itemsep=4pt]
  \item $-1$ and $-\dfrac12$.
  \item $225^\circ$.
  \item $x=\dfrac\pi3$ and $x=\dfrac{2\pi}{3}$.
  \item $x=\dfrac{3\pi}{4}$ and $x=\dfrac{7\pi}{4}$.
  \item $\cos x=\pm\dfrac{\sqrt3}{2}$: $x=\dfrac\pi6,\ \dfrac{5\pi}{6},\ \dfrac{7\pi}{6},\ \dfrac{11\pi}{6}$.
  \item $\sin^2\theta=1-\dfrac49=\dfrac59$, and sine is positive in quadrant II: $\sin\theta=\dfrac{\sqrt5}{3}$, $\tan\theta=-\dfrac{\sqrt5}{2}$.
  \item $\dfrac{\sin x}{\cos x}+\dfrac{\cos x}{\sin x}=\dfrac{\sin^2x+\cos^2x}{\sin x\cos x}=\dfrac{1}{\sin x\cos x}=\sec x\csc x$.
  \item $\dfrac\pi3$, \ $-\dfrac\pi3$ \ and \ $\dfrac\pi2$.
  \item \begin{enumerate}[label=(\alph*),leftmargin=*,itemsep=3pt]
    \item $\cos\dfrac{2\pi}{3}+3\cos\dfrac\pi3+2=-\dfrac12+\dfrac32+2=3$.
    \item $2\cos^2x-1+3\cos x+2=2\cos^2x+3\cos x+1$.
    \item $(2\cos x+1)(\cos x+1)=0$. $\cos x=-\dfrac12$ gives $\dfrac{2\pi}{3}$ and $\dfrac{4\pi}{3}$; $\cos x=-1$ gives $\pi$.
  \end{enumerate}
\end{enumerate}
\end{solbox}

\end{document}
